% Chapter 5: System architecture and autonomy. First draft of the autonomy mathematics.
% Every equation and default value below is transcribed from the Terra repository at
% main 293acc1 (Oct 24, 2026):
%   crates/terra-state/src/lib.rs            (velocity estimation)
%   crates/terra-mapping/src/lib.rs          (occupancy update)
%   crates/terra-control/src/{lib,turning}.rs (velocity controller, turning policy)
%   crates/terra-navigation/src/{lib,astar}.rs (local planner, A*)
%   crates/terra-waypoint/src/lib.rs         (pursuit, local projection)
%   crates/terra-autonomy/src/arbiter/motion.rs (per-level composition)
%   crates/terra-autonomy/src/search.rs      (target confirmation)
%   crates/terra-exploration/src/*.rs        (frontier exploration)
%   crates/terra-autonomy/src/outdoor.rs, outdoor/heading.rs (GPS guidance, heading fusion)
%   crates/terra-types/src/timing.rs         (time constants)
% plus Terra PR #60 and ARGOS PR #64 (merged Oct 25, 2026) for L4 Explore.
% Updated Oct 26, 2026 against Terra main 0110519 and ARGOS main 552c1e7:
%   crates/terra-runtime/src/ground.rs        (level ground plane, Terra #66; mount floor, #68)
%   crates/terra-types point_capture, terra-mapping label_depth (Terra #64, #67)
%   crates/terra-mobile/src/actuators.rs      (Rust actuator routing in shadow mode, Terra #63)
%   crates/terra-autonomy/src/task.rs, arbiter/task.rs, docs/autonomy/TASK.md (L5, Terra #70)
%   ARGOS #69 (focused subscriptions), #72/#83 (Capture 3D, Live 3D), #85 (Task L5 panel).
% Updated Oct 27, 2026 against Terra main 4d63878 and ARGOS main (after #89):
%   Terra #71/#72 (rover profiles, wheel deadband), #74 (IMU through the mount),
%   #75 (L4 look-around cleared on level change), #76/#77/#78 (ARKit mesh floor veto,
%   floor walk off the control queue), #80/#81 (3D classification page), #82 (sensor
%   hold causes, planner window, autonomy brief), #83 (audit fixes); ARGOS #87, #89.
%   Terra #84 / ARGOS #90 (barrier label for glass), merged Oct 27, untested on real glass.
%   Velocity validation: OptiTrack sessions of Oct 26--27, data and scripts in the
%   project's velocity-experiment folder (analyze_velocity.py, mocap_session).
% Defaults are software starting values, not calibrated physical limits.
% Section sec:auto-fusion is transcribed from Terra scripts/plot_sim_rovers.py
% (uncommitted working copy, Oct 26, 2026) and its Oct 25 technical report.
% Section sec:arch is a first outline; the topic contract and logging still need detail.

\chapter{System Architecture and Autonomy}
\label{ch:autonomy}

This chapter describes the software that runs the rover and the mathematics behind each autonomy level.
Section~\ref{sec:arch} outlines the architecture.
The rest of the chapter gives the models and control laws in the order a command passes through them: state estimation and mapping (Sections~\ref{sec:auto-state}--\ref{sec:auto-map}), the vehicle model and velocity controller (Section~\ref{sec:auto-control}), the turning policy (Section~\ref{sec:auto-turning}), the local planner (Section~\ref{sec:auto-planner}), and then the intent each level produces (Section~\ref{sec:auto-levels}).
Section~\ref{sec:auto-outdoor} covers outdoor GPS guidance and heading fusion, and Section~\ref{sec:auto-fusion} a prototype that fuses two rovers' maps.

Every equation in this chapter is taken from the Terra source code, and every default value is the one in the code as of October 2026.
They are software starting points, chosen in simulation; Chapter~\ref{ch:rover} lists the physical limits that still have to be measured.

\section{Architecture}
\label{sec:arch}

The rover software has two layers on the phone.
TerraPhone is a Swift application.
It owns everything that needs an Apple framework: the ARKit session and depth frames, Core Motion, GPS and compass, the Bluetooth link to the motor bridge, WebRTC first-person video, and, for L5, the on-device language model (Section~\ref{sec:auto-l5}).
Below it is the Terra core, a set of Rust crates compiled for iOS and called through UniFFI bindings.
The same crates are linked by the Zorvane simulator (Chapter~\ref{ch:simulation}), so the autonomy that runs on the phone is the autonomy that runs in simulation.
An Android client for the same Bluetooth actuator protocol also exists, but the phone used in this thesis is an iPhone.

\subsection{Where the control loop belongs}
\label{sec:arch-boundary}

The project follows one rule for splitting work between the two languages: Swift gathers sensor data, talks to the hardware and the network services that only Apple's frameworks provide, and builds the occupancy input from depth; the control loop, safety gating, mode and tick policy, and actuator routing belong in Rust.
Rust code is shared with the simulator and covered by its tests, so a rule that lives there is exercised in both places; a rule in Swift runs only on a phone.

In October 2026 the code is part of the way there.
The arbiter, planners, velocity controller and level logic are in Rust, but TerraPhone still ticks the loop and still holds the motor gates and the mapping from wheel effort to actuator outputs.
The first migration step (Terra PR~\#63) makes Rust route each tick's efforts through the rover's drive profile as well and return the resulting actuator frame; TerraPhone still sends its own frame and logs any tick where the two differ.
This \emph{shadow mode} lets the Rust routing be checked against the Swift routing on real drives before it takes over.
\todo{Report how many shadow-mode mismatches were logged on real drives, then switch the actuator output to the Rust frame and describe the final split.}

Two fixes from the same period show why the loop's timing is delicate on a phone.
ARKit delivered frames on the same queue as the 100~Hz control tick, depth mapping and point-cloud publishing, so frames that arrived during a step waited holding their \texttt{ARFrame}, and ARKit warned that eleven frames were retained (fixed in Terra PR~\#65).
A semaphore wait between the high-priority frame thread and a utility-priority capture queue was a priority inversion (fixed in Terra PR~\#61).
Neither changes the mathematics below, but both affect how fresh the inputs to it are.

An audit of TerraPhone and its Rust bindings on October 27, 2026 measured how far the split still has to go.
The Rust \texttt{step} already returns the routed actuator frame, but iOS only compares it with its own and still encodes and sends the Swift frame; the Android client already sends the Rust one.
The e-stop latch, the mount-confirmation gate, the pause on lost visual tracking and the tick dividers exist only in Swift, and iOS and Android use different conventions for the phone-to-body mount transform.
The audit also found bugs in the Swift layer, the most serious being that stopping the controller cleared the e-stop latch, so a stop followed by a restart silently re-armed a rover that had been e-stopped.
Those bugs and the dead code around them were fixed in Terra PR~\#83: the latch now survives a stop and is cleared only by an explicit reset.
The remaining order of work is to send the Rust actuator frame on iOS, give the Rust tick a typed result in place of a status string, and move the gates and tick dividers into Rust.
\todo{Report the migration once iOS sends the Rust actuator frame.}

\subsection{Loop timing on the phone}
\label{sec:arch-timing}

Every stage of the pipeline in this chapter holds its output when an input is stale: the estimator when the IMU is older than 0.1~s, the controller when its time step exceeds 0.1~s, the turning limiter when its step exceeds 0.1~s.
These rules make the rover stop rather than act on old data, but they also turn any delay on the control queue into a hold, and a hold clears the controller's integrators.
Three faults found during the motion-capture sessions of October 26--27 (Section~\ref{sec:auto-state-validation}) were of this kind, and none was a fault in the control law itself.

\begin{description}
\item[Mesh walk on the control queue.] The floor veto of Section~\ref{sec:auto-floor-veto} first walked the ARKit mesh on the control queue once a second. Each walk stalled pose and IMU delivery for about 150~ms, so the estimator reported itself not ready, the controller held, and a commanded 0.4 in effort fell to about 0.19 after each reset. The walk now runs on a utility queue and is built outside the map lock (Terra PR~\#77), and walks in flight are dropped after a map or ground reset (Terra PR~\#78).
\item[Status parsing at L3.] In L3 the autonomy status carries the planner trace, about 30~KB of JSON, and TerraPhone parsed it three to four times on every 10~ms tick. The control queue fell behind, the IMU sample went stale, and the rover held with \texttt{sensor\_unhealthy}; one L3 run logged 191 such holds. Rust now also returns a short autonomy brief of under 1~KB that the hot Swift paths read instead (Terra PR~\#82). Before the fix, ARGOS could only show that the rover was holding; now Terra records the cause of each sensor hold and keeps it for 10~s, and ARGOS shows it on a \emph{Sensor hold} row (Terra PR~\#82, ARGOS PR~\#89). The cause read \texttt{stale\_imu}, which is what located the fault.
\item[Late ticks in straight runs.] In the fourth motion-capture session, five straight runs ended about one second after they started, when one tick arrived more than 0.1~s late and the turning limiter held. The teleoperation client then treated the hold as a release of the drive keys. This is probably the same status-parsing load, since that phone build did not yet have Terra PR~\#82, but it has not been confirmed.
\end{description}

The holds explain why L3 failed to reach a goal that L2 reached from the same start.
L2 bypasses the planner, while the L3 planner's velocity window was centred on the measured speed (Section~\ref{sec:auto-planner}): after each hold the measured speed of a wheel inside the motor deadband was zero, the planner offered at most 0.1~m/s, and the wheels never left the deadband.
\todo{Measure the control-tick latency distribution on the phone at L0 and L3 after Terra PR~\#82, and report how often a tick exceeds the 0.1~s limits.}

\subsection{Fleet communication}
\label{sec:arch-topics}

Every rover publishes and subscribes under its own numeric identifier below one shared Zenoh prefix (Chapter~\ref{ch:competition}).
Light telemetry, such as pose, localization, goal and autonomy status, hardware state, and outdoor navigation, is received for every rover.
Heavy topics, the occupancy grid, the point cloud and depth frames, are subscribed only for the rover the operator has selected (ARGOS PR~\#69), so the router does not forward every rover's map to the operator station.
Team CERBERUS made the same choice in the DARPA Subterranean Challenge, where camera streams were sent only for the robot the Human Supervisor selected~\citep{tranzatto2022finals}.
Requests that are not motion, such as a point-cloud capture request, go to their own inbox outside the motion-authority path, and bulk replies are sent at background priority and dropped rather than queued when the link is congested.
\todo{Describe run identity and authority revisions, the Pi Bluetooth protocol and logging in an appendix, and list the topic contract.}

\subsection{Motion authority}

For the rest of this chapter one property of the architecture matters: all motion authority sits in one place, the \emph{arbiter} in the Terra core on the phone.
Operator commands, waypoint goals, and autonomous goals are all \emph{intents}.
The arbiter turns the intent of the active level into a velocity command, checks it, and hands it to the velocity controller.
The control loop ticks every 10~ms; wheel commands are sent to the motor bridge every fifth tick, at 20~Hz; the local planner works on a 0.1~s step.

\section{Notation and frames}
\label{sec:auto-notation}

The rover pose in the local map frame is $\mathbf{p} = (x, y, \theta)$, with $x, y$ in metres and yaw $\theta$ in radians, measured counter-clockwise from the $+x$ axis.
The body twist is $(v, \omega)$: forward speed in m/s and yaw rate in rad/s, positive to the left.
A hat marks an estimate ($\hat v$), a star a target ($v^\ast$), and a subscript $i$ the intent produced by an autonomy level ($v_i, \omega_i$).
$\operatorname{clip}(a, \pm b)$ limits $a$ to $[-b, b]$, and $\operatorname{wrap}(\cdot)$ maps an angle to $[-\pi, \pi)$.
Geographic headings are in degrees, clockwise from true north.

\section{State estimation}
\label{sec:auto-state}

The velocity controller needs the rover's forward speed and yaw rate.
The phone has no wheel encoders in the loop, so Terra estimates them from ARKit's visual--inertial odometry (VIO) and the Core Motion IMU.
An independent benchmark found ARKit the most stable and accurate of four commercial VIO systems~\citep{kim2022benchmark}, so its pose is a reasonable stand-in for wheel odometry.
The latest VIO sample, at time $t_0$, gives an orientation quaternion $\mathbf{q}_0$ and a world-frame velocity $\mathbf{v}_0$.
Because VIO arrives late and at a lower rate than the IMU, the estimator replays the retained IMU samples $(\mathbf{a}_k, \boldsymbol{\omega}_k)$ from $t_0$ to the current time.
For each interval $\delta_k$,
\begin{align}
  \mathbf{q}_{k+\frac12} &= \mathbf{q}_k \otimes \exp\!\left(\tfrac12 \boldsymbol{\omega}_k \delta_k\right), \\
  \mathbf{v}_{k+1} &= \mathbf{v}_k + R(\mathbf{q}_{k+\frac12})\,\mathbf{a}_k\,\delta_k, \\
  \mathbf{q}_{k+1} &= \mathbf{q}_k \otimes \exp\!\left(\boldsymbol{\omega}_k \delta_k\right),
\end{align}
where $\exp(\cdot)$ maps a rotation vector to a unit quaternion and $R(\mathbf{q})$ is the corresponding rotation.
The estimate is the body-frame forward component and the latest gyro yaw rate:
\begin{equation}
  \hat v = \left[R(\mathbf{q}_N)^{\top} \mathbf{v}_N\right]_x, \qquad \hat\omega = \omega_{z,N}.
\end{equation}
The estimate is withheld, and the controller stops, if VIO reports lost tracking, if the IMU is older than 0.1~s, or if VIO is older than 0.35~s.
On the phone, the VIO velocity $\mathbf{v}_0$ is itself a finite difference of the ARKit position of the rover's reference point, low-pass filtered with a 30~ms time constant.

The IMU vectors must be rotated from the phone's axes into the rover body through the mount calibration (Chapter~\ref{ch:rover}).
Until Terra PR~\#74 the iOS IMU path assumed a phone lying flat, so on an upright mount the gyro axis read as yaw was the phone's roll axis, and the estimated yaw rate stayed near zero through arcs at 0.5 to 1.1~rad/s.
The fix maps Core Motion through the saved mount transform.
That transform is only as good as the saved mount: with the landscape phone on CrabTrapBot, a calibration with roll 0 leaves the yaw-rate axis and body up 90$^\circ$ off.
The ARKit heading and forward speed do not depend on this mapping and were unaffected.

\subsection{Validation against motion capture}
\label{sec:auto-state-validation}

Neither chassis has wheel encoders, so the VIO speed is the only speed the controller sees.
To check it, CrabTrapBot was driven in an OptiTrack motion-capture volume of about $12 \times 8$~m on October 26 and 27, 2026.
OptiTrack (Motive~3, NatNet~4) tracked a rigid body on the chassis at 120~Hz; a 3~m tape measure read 0.01~m different from OptiTrack, so it served as ground truth.
A bridge script republished the rigid-body pose on the fleet router, and a single program on the laptop recorded both poses and drove the rover in straight lines at a set speed and distance through ARGOS-style teleoperation.
The phone's observability journal was copied off the phone after each session.
For every run, the steady part of the phone's forward speed, after the speed-up and before the stop, was compared with OptiTrack's.

Table~\ref{tab:vio-scale} gives the ratio of phone to OptiTrack speed for the 17 usable straight runs at 0.2 to 0.4~m/s.

\begin{table}[ht]
\centering
\caption{Steady forward speed, phone VIO divided by OptiTrack, for straight runs on CrabTrapBot, October 26--27, 2026.}
\label{tab:vio-scale}
\begin{tabular}{@{}l r l@{}}
\toprule
Session & Runs & Ratio per run \\
\midrule
1 (Oct 26) & 2 & 0.98, 0.81 \\
2 (Oct 26) & 4 & 0.81, 0.89, 0.93, 0.88 \\
3 (Oct 26) & 1 & 1.01 \\
4 (Oct 27) & 10 & 0.83, 0.60, 0.93, 0.95, 0.92, 0.86, 0.90, 1.01, 0.92, 0.97 \\
\midrule
All & 17 & mean 0.894, s.d.\ 0.098, median 0.92 \\
Without the 0.60 run & 16 & mean 0.912, s.d.\ 0.065 \\
\bottomrule
\end{tabular}
\end{table}

The phone under-read the speed in 15 of the 17 runs.
The same under-read appeared when the phone was carried by hand along a line, 4.06~m on the phone against 4.42~m on OptiTrack (0.92), so it comes from ARKit and not from the wheels or motors.
The other properties the controller needs were good: at rest the speed bias was 0.0009~m/s with a standard deviation of 0.003~m/s; while driving the noise was 0.007~m/s in the control band and 0.012~m/s above it; and ARKit's heading agreed with OptiTrack to within 0.1$^\circ$ over a 475$^\circ$ arc.
The latency relative to the gyro was roughly 50 to 70~ms, a rough figure because the journal logged at only 8~Hz.
Two events fell outside this behaviour: the 0.60 run and a 2.75~m jump in ARKit position on one arc, both consistent with VIO losing and re-finding its scale or anchor.

For the controller of Section~\ref{sec:auto-control}, a scale error of about 9\% shifts the loop gain by the same small amount and leaves it stable, but the loop drives the \emph{measured} speed to the target, so the rover really travels about 9\% faster and farther than it believes.
Over a 1~s planner rollout at 0.3~m/s that is 3 to 4~cm, less than one grid cell.
The occupancy map is built from LiDAR depth in the same ARKit frame, so the rover stays consistent with its own map; absolute distances, GPS goals and the fusion of two rovers' maps are where the error shows.
The leading explanation is the room itself: OptiTrack's infrared strobes may interfere with the phone's LiDAR, and the floor was plain.
\todo{Repeat the straight runs in a textured hallway and outdoors with a tape measure; if the ratio stays near 0.91, add a fixed scale correction to the velocity estimate.}
\todo{Measure arcs and turns after setting the landscape mount roll, log at 100~Hz to tighten the latency estimate, and sweep 0.1 to 0.5~m/s.}
\todo{Add a guard that rejects VIO position steps larger than the rover could physically move, and size it from longer sessions.}

\section{Occupancy mapping}
\label{sec:auto-map}

The local map is a rolling grid of $200 \times 200$ cells at resolution $r = 0.1$~m.
Each cell holds a log-odds value $\ell$.
For each depth frame, every sampled pixel $(u, v)$ with depth $z$ is back-projected through the camera intrinsics,
\begin{equation}
  \mathbf{r}_c = \left(\frac{(u - c_x) z}{f_x},\; \frac{(v - c_y) z}{f_y},\; z\right),
\end{equation}
rotated into the world frame, and classified as an obstacle hit if its endpoint lies between 0.15 and 2.0~m above the ground plane, within 15~m of the camera.
Each pixel gets one of six labels: free, obstacle, overhead, out of range, outside the band, or no return, plus the floor-veto label of Section~\ref{sec:auto-floor-veto}; the grid update below uses the same classifier, so the labels and the map evidence cannot disagree.
The cells along the ray, found by Bresenham traversal, are marked free, and the endpoint cell is marked occupied for a hit.
Each cell is updated at most once per frame, with obstacle evidence taking precedence:
\begin{equation}
  \ell \leftarrow \operatorname{clip}\!\left(\ell + \begin{cases} +0.85 & \text{occupied} \\ -0.40 & \text{free} \end{cases},\; \pm 4\right).
\end{equation}
The map published to the planner stores $P = 100\,\sigma(\ell) = 100 / (1 + e^{-\ell})$ for observed cells and $-1$ for unobserved ones.
A cell is \emph{occupied} if $P \ge 65$, that is if $\ell \ge \ln(65/35) \approx 0.62$, \emph{free} if $0 \le P < 65$, and \emph{unknown} otherwise.
The asymmetric increments mean that one hit outweighs two misses, so a thin obstacle seen intermittently stays in the map.

\subsection{The ground plane}
\label{sec:auto-ground}

Every height in the classifier is measured from a ground plane, so an error in the plane becomes an error in the map.
The plane is level, with normal along world up, and lies a vertical distance $h_c$ below the camera:
\begin{equation}
  z_{\text{ground}} = z_{\text{cam}} - h_c,
\end{equation}
so it follows the terrain's elevation as the rover moves, but never its tilt.
The body's up axis is used only to reject a tipped-over rover.
An earlier version used the body's up axis as the plane normal.
A labelled point capture on a flat floor showed the body axis about 16$^\circ$ off gravity because of mount pitch and roll, so the plane cut through the floor: on one side, floor beyond about 2~m was labelled obstacle, and on the other, walls up to 2.5~m high read as free (fixed in Terra PR~\#66).

The camera height $h_c$ comes from the mount calibration (Chapter~\ref{ch:rover}) when one with a plausible height (0.1 to 3~m) is saved, as a robot model would use a fixed camera transform (Terra PR~\#68).
Without a calibration, $h_c$ is taken once from the first ARKit horizontal plane of at least 0.4~m$^2$ whose height stays within 3~cm for 1~s.
That fallback is fragile: if the app is started with the phone in hand and the phone is mounted afterwards, $h_c$ is the hand height and the plane sinks below the real floor.

Before a floor is found, the map uses a \emph{horizon band}: returns within $\pm10^\circ$ of the horizon count as obstacles out to the range at which the band meets the floor, $h_c/\tan 10^\circ$, about 1.3~m for a 0.23~m camera.
The band has no margin, so a camera height set a little too high puts that limit beyond the true floor intercept: in one capture in the motion-capture room every one of the 4,094 obstacle points was floor, in a ring 1.1 to 1.3~m around the rover.
The band also persists if the ground plane rejects the body's up axis, which it does when the body appears tilted by more than 60$^\circ$.
On CrabTrapBot that happened because the phone is mounted in landscape and the mount calibration still had roll 0, so the body appeared to lie on its side.
The map status now reports the tilt, and band frames now use the same height offset as the mesh floor, which had been drawn 0.23~m too high (Terra PR~\#82).
\todo{Repeat the flat-floor point capture with the level, mount-calibrated plane on CrabTrapBot and report the fraction of floor points labelled obstacle before and after.}

\subsection{Inspecting the classifier}
\label{sec:auto-capture}

The ground-plane fault was found with a diagnostic built for that purpose.
On request from ARGOS, the rover sends every depth pixel, projected to world coordinates and labelled by the rules above, as binary frames (Terra PR~\#64, ARGOS PR~\#72).
A \emph{Capture 3D} records 10~s at 1~Hz, about 640~KB per frame, and ARGOS saves it as a PLY file coloured by label.
A \emph{Live 3D} stream sends decimated frames at up to 5~Hz for as long as ARGOS renews a lease, and ARGOS draws the newest frame on the map under the rover (Terra PR~\#67, ARGOS PR~\#83); ARGOS uses a 6~s lease renewed every 2~s, 2~Hz, and every third pixel in each direction, about 70~KB per frame.
Decimation scales the camera intrinsics so that every ray kept is unchanged.
The always-on thinned point cloud that the phone used to stream is now sent only while a capture runs.
For work without a rover or a laptop, TerraPhone also has a \emph{3D classification} debug page that shows the live labelled point cloud and the mesh floor cells on the phone itself, held in the hand (Terra PRs~\#80 and \#81).

\subsection{Floor veto from the ARKit mesh}
\label{sec:auto-floor-veto}

A level ground plane removes the tilt error but not every false obstacle: depth noise, a slightly wrong camera height or a gentle slope can still lift a floor return above the 0.15~m band floor.
ARKit can classify its reconstructed scene mesh into surfaces such as floor, wall, window and door, and this gives a second, independent estimate of where the floor is.
TerraPhone turns on the classified mesh and sends the centres of floor faces to the Rust map once a second.
The map bins them into 0.1~m cells, keeps the median height per cell, and drops cells more than 8~cm from the median of the whole surface.
In the classifier, a return that would be an obstacle but lies less than 0.15~m above a trusted mesh floor cell is relabelled \emph{floor veto} and clears its cell like a free return; a return higher than that above the same cell is still an obstacle, so an object standing on the floor is not erased.
The mesh floor also seeds the ground-plane height when no floor plane anchor has been found.
Captures carry the same surface, so a captured frame shows exactly what the map did, and ARGOS colours the new label (Terra PR~\#76, ARGOS PR~\#87, both merged October 26, 2026).
Two design questions remain open: whether a veto should count as free evidence or as no evidence, and how well ARKit's floor classification holds on turf.
\todo{Repeat the flat-floor capture with the floor veto on and report the fraction of floor returns labelled obstacle with and without it.}

\subsection{Glass}
\label{sec:auto-glass}

The phone's LiDAR sees through clear glass, so the map carves glass walls and doors as free space, and L3 would route through them.
A change merged on October 27, 2026 uses the same classified mesh: faces ARKit labels wall, window or door are sent to the map, and cells with at least three such faces are marked as a \emph{barrier} with a height span.
A depth ray that crosses a barrier cell stops there; the cell is marked occupied and the return gets a new barrier label (Terra PR~\#84, ARGOS PR~\#90).
This only works where ARKit has put mesh faces on the glass, which a clear pane seen at a shallow angle may not get.
\todo{Test the barrier on real glass doors and windows; if it misses them, compare remembering bumps and adding an ultrasonic sensor.}

\section{Vehicle model and velocity control}
\label{sec:auto-control}

\subsection{Kinematics}

Both chassis are differential drive.
Treated as a unicycle,
\begin{equation}
  \dot x = v \cos\theta, \qquad \dot y = v \sin\theta, \qquad \dot\theta = \omega,
  \label{eq:unicycle}
\end{equation}
with the wheel speeds related to the body twist by
\begin{equation}
  v = \frac{v_R + v_L}{2}, \qquad \omega = \frac{v_R - v_L}{b_{\text{eff}}},
\end{equation}
where $b_{\text{eff}}$ is the effective track width.
On a skid-steer chassis $b_{\text{eff}}$ is larger than the geometric track because the wheels scrub sideways in a turn.
The controller below never uses $b_{\text{eff}}$ explicitly: it works in normalized wheel effort, and the track width, motor gain, and scrub are absorbed into its feedforward and feedback gains.
Equation~\eqref{eq:unicycle} is the model the planner uses to predict trajectories.

\subsection{Velocity controller}

The velocity controller turns a target twist $(v^\ast, \omega^\ast)$ into normalized left and right wheel efforts $u_L, u_R \in [-1, 1]$.
Targets are first limited to $|v^\ast| \le 2$~m/s and $|\omega^\ast| \le 2$~rad/s.
With errors $e_v = v^\ast - \hat v$ and $e_\omega = \omega^\ast - \hat\omega$, the linear and yaw efforts are feedforward plus proportional--integral terms,
\begin{align}
  u_v &= k^{\text{ff}}_v v^\ast + k^{p}_v e_v + I_v, \label{eq:uv}\\
  u_\omega &= k^{\text{ff}}_\omega \omega^\ast + k^{p}_\omega e_\omega + I_\omega. \label{eq:uw}
\end{align}
They are mixed into wheel efforts and scaled down together, which preserves the ratio between turning and driving when the command saturates:
\begin{equation}
  s = \max\!\left(1,\; \frac{\max(|u_v - u_\omega|,\, |u_v + u_\omega|)}{u_{\max}}\right), \qquad
  u_L = \frac{u_v - u_\omega}{s}, \quad u_R = \frac{u_v + u_\omega}{s}.
  \label{eq:mix}
\end{equation}
Unless turning in place is allowed, the wheels may not counter-rotate: if $u_L u_R < 0$, the wheel opposing the mean direction $\tfrac12(u_L + u_R)$ is set to zero.
The integrators use back-calculation anti-windup against the effort actually applied, $\bar u_v = \tfrac12(u_L + u_R)$ and $\bar u_\omega = \tfrac12(u_R - u_L)$:
\begin{align}
  I_v &\leftarrow \operatorname{clip}\!\left(I_v + \left[k^{i}_v e_v + k_{\text{aw}}(\bar u_v - u_v)\right]\Delta t,\; \pm u_{\max}\right), \\
  I_\omega &\leftarrow \operatorname{clip}\!\left(I_\omega + \left[k^{i}_\omega e_\omega + k_{\text{aw}}(\bar u_\omega - u_\omega)\right]\Delta t,\; \pm u_{\max}\right).
\end{align}
When the target is zero and the rover is nearly still ($|\hat v|, |\hat\omega| < 0.015$), the outputs are forced to zero and the integrators reset, so the motors do not hum at rest.
If in-place turning is allowed and the target is a pure rotation, the controller bypasses the loop and applies a fixed spin effort, $u_R = -u_L = u_s \operatorname{sgn}\omega^\ast$.
Small DC motors do not turn below some effort $d$, so a rover profile may set a wheel deadband (Terra PR~\#71).
Each moving wheel's effort $u$ is then lifted past it before it is sent,
\begin{equation}
  u' = \operatorname{sgn}(u)\left(d + (u_{\max} - d)\,\frac{|u|}{u_{\max}}\right) \quad \text{for } |u| \ge 0.01,
\end{equation}
while efforts under 0.01 stay zero and the integrators work on the uncompensated effort.
The default is $d = 0$; Chapter~\ref{ch:rover} gives the per-rover values.

The controller stops (zero effort, integrators cleared) when the target is 0.5~s old or older, when the time step is non-positive or longer than 0.1~s, or when the estimator is not healthy.
Table~\ref{tab:controller-gains} lists the defaults.
With $k^{\text{ff}}_v = 1/3$, full effort corresponds to a feedforward speed of 3~m/s.

\begin{table}[ht]
\centering
\caption{Velocity controller defaults (\texttt{ControllerConfig}). \todo{tune on CrabTrapBot and report the tuned values}}
\label{tab:controller-gains}
\begin{tabular}{@{}l l r@{}}
\toprule
Symbol & Meaning & Default \\
\midrule
$k^{p}_v,\; k^{i}_v$ & Linear proportional and integral gains & 0.3,\; 0.6 \\
$k^{p}_\omega,\; k^{i}_\omega$ & Yaw proportional and integral gains & 0.1,\; 0.3 \\
$k^{\text{ff}}_v,\; k^{\text{ff}}_\omega$ & Linear and yaw feedforward & $1/3$,\; 0.2 \\
$k_{\text{aw}}$ & Anti-windup gain & 4.0 \\
$u_{\max}$ & Maximum effort & 1.0 \\
$u_s$ & In-place spin effort & 0.35 \\
$d$ & Wheel deadband & 0 \\
-- & Target timeout; maximum step & 0.5~s;\; 0.1~s \\
\bottomrule
\end{tabular}
\end{table}

\section{Turning policy}
\label{sec:auto-turning}

Skid-steer turns scrub the tyres and disturb the phone's visual tracking, so every \emph{managed} command, meaning ARGOS teleoperation and every autonomous level, passes through a turning limiter before the controller.
Local joystick driving on the phone does not.
Given a requested $\omega^\ast$ at forward speed $v$, the limiter first caps the yaw rate by a minimum turning radius $R_{\min}$,
\begin{equation}
  \omega_{\text{cap}} = \begin{cases} \omega_{\max} & \text{in-place turns allowed,} \\ \min(\omega_{\max},\; |v| / R_{\min}) & \text{otherwise,} \end{cases}
\end{equation}
so that the rover cannot turn on the spot by default.
It then slews from the previous command $\omega_{\text{prev}}$ toward the capped request $\omega_d = \operatorname{clip}(\omega^\ast, \pm\omega_{\text{cap}})$, passing through zero on a reversal:
\begin{align}
  g &= \begin{cases} 0 & \omega_{\text{prev}}\,\omega_d < 0, \\ \omega_d & \text{otherwise,} \end{cases} \qquad
  \alpha = \begin{cases} \alpha_{\text{dec}} & \omega_{\text{prev}}\,\omega_d < 0 \text{ or } |g| < |\omega_{\text{prev}}|, \\ \alpha_{\text{acc}} & \text{otherwise,} \end{cases} \\
  \omega &= \operatorname{clip}\!\left(\omega_{\text{prev}} + \operatorname{clip}(g - \omega_{\text{prev}},\, \pm\alpha\Delta t),\; \pm\omega_{\text{cap}}\right).
\end{align}
The defaults are $\omega_{\max} = 0.5$~rad/s, $\alpha_{\text{acc}} = 0.5$~rad/s$^2$, $\alpha_{\text{dec}} = 0.8$~rad/s$^2$, $R_{\min} = 0.5$~m, and in-place turning off; reverse speed is limited to 0.15~m/s.
When an autonomous level needs to face a new direction from rest, the arbiter plans a multi-point manoeuvre of forward and reverse arcs that respects $R_{\min}$, rather than spinning.

\section{Local planner}
\label{sec:auto-planner}

The local planner is the safety layer for every level that checks obstacles.
It is a deterministic velocity-space sampler in the family of the dynamic window approach~\citep{fox1997dwa}: it searches the velocities reachable in one step for the one closest to the intent whose predicted path stays clear.

\subsection{Dynamic window}

The window grows from a centre twist $(\bar v, \bar\omega)$.
For each axis the centre is the last command the arbiter allowed, $(v_{\text{prev}}, \omega_{\text{prev}})$, zero after a hold, unless the rover is already measured going faster in the same direction:
\begin{equation}
  \bar v = \begin{cases} \hat v & v_{\text{prev}}\,\hat v \ge 0 \text{ and } |\hat v| > |v_{\text{prev}}|, \\ v_{\text{prev}} & \text{otherwise,} \end{cases}
\end{equation}
and likewise for $\bar\omega$.
With planning step $\Delta t = 0.1$~s, the reachable window is
\begin{align}
  V &= [\bar v - a\Delta t,\; \bar v + a\Delta t] \cap [-v_{\text{rev}},\; v_{\max}], \\
  \Omega &= [\bar\omega - \alpha\Delta t,\; \bar\omega + \alpha\Delta t] \cap [-\omega_{\max},\; \omega_{\max}].
\end{align}
Until Terra PR~\#82 the window was centred on the measured twist alone.
A wheel inside the motor deadband measures zero, so after a hold the planner offered at most $a\Delta t = 0.1$~m/s, which kept the wheels in the deadband; this is how L3 stalled where L2 did not (Section~\ref{sec:arch-timing}).
The planner samples an $11 \times 11$ grid of $(v, \omega) \in V \times \Omega$, evenly spaced including the endpoints.
For managed commands each sample is first passed through the turning limiter of Section~\ref{sec:auto-turning}, so the planner only considers commands that will actually be executed.

\subsection{Trajectory rollout with braking}

Each candidate is rolled forward under Equation~\eqref{eq:unicycle} for the prediction horizon $H$, followed by a braking tail long enough to stop at deceleration $a_b$:
\begin{equation}
  T = H + \frac{|v|}{a_b}, \qquad
  \lambda(t) = \begin{cases} 1 & t \le H, \\ \max\!\left(0,\; 1 - \dfrac{(t - H)\,a_b}{|v|}\right) & t > H. \end{cases}
\end{equation}
Both velocities are scaled by $\lambda(t)$, which preserves the path's curvature while braking.
With step $\delta = \min\!\left(r / (2\max(|v|, 0.1)),\; 0.1\right)$, so that the rover advances at most half a cell per step,
\begin{equation}
  \theta_{n+1} = \theta_n + \lambda\omega\delta, \quad
  x_{n+1} = x_n + \lambda v \cos\theta_{n+1}\,\delta, \quad
  y_{n+1} = y_n + \lambda v \sin\theta_{n+1}\,\delta.
\end{equation}

\subsection{Admissibility}

The rover footprint is a disc of radius $\rho$.
A pose is clear if no cell whose centre lies within $\rho + 0.7072\,r$ of the rover centre is occupied, outside the map, or unknown; the extra half cell diagonal makes the test conservative.
Unknown cells are tolerated only under the rover's footprint at its starting pose, so that it can drive off its own unobserved underbody, and never when reversing.
A candidate is \emph{admissible} if every pose of its rollout, including the braking tail, is clear.
Because the tail is included, an admissible command is one the rover can still stop from before reaching anything the map knows about, or anything it has not seen.

\subsection{Selection}

Among admissible candidates the planner minimizes
\begin{equation}
  J(v, \omega) = (v - v_i)^2 + 0.3\,(\omega - \omega_i)^2 + \lambda_{\text{rev}}\max(0, -v),
  \label{eq:planner-cost}
\end{equation}
with ties broken by sample order.
If no candidate is admissible the command is zero and the arbiter reports \texttt{obstacle\_blocked}.
If the chosen command differs from the intent by more than 0.15 in either component the decision is labelled \texttt{assisted}, which is how the operator station shows that the planner, not the operator or the follower, is shaping the motion.
Table~\ref{tab:planner} lists the defaults.

\begin{table}[ht]
\centering
\caption{Local planner defaults (\texttt{PlannerConfig}).}
\label{tab:planner}
\begin{tabular}{@{}l l r@{}}
\toprule
Symbol & Meaning & Default \\
\midrule
$\rho$ & Footprint radius & 0.65~m \\
$v_{\max},\; v_{\text{rev}}$ & Forward and reverse speed limits & 2.0,\; 0.15~m/s \\
$\omega_{\max}$ & Yaw rate limit & 2.0~rad/s \\
$a,\; \alpha$ & Linear and yaw acceleration & 1.0~m/s$^2$,\; 2.0~rad/s$^2$ \\
$a_b$ & Braking deceleration & 1.0~m/s$^2$ \\
$H$ & Prediction horizon & 2.0~s (1.0~s in L4 Explore) \\
$\lambda_{\text{rev}}$ & Reverse penalty & 0.05 \\
$\Delta t$ & Planning step & 0.1~s \\
\bottomrule
\end{tabular}
\end{table}

\section{Control laws for each autonomy level}
\label{sec:auto-levels}

Each level differs only in how it produces the intent $(v_i, \omega_i)$ and in which of the layers above it uses.
Figure~\ref{fig:pipeline} shows the pipeline, and Table~\ref{tab:level-math} summarizes it.

\begin{figure}[ht]
\centering
\begin{tikzpicture}[
  box/.style={draw, rounded corners, align=center, minimum width=2.0cm, minimum height=0.9cm, font=\small},
  lbl/.style={font=\scriptsize, align=center},
  >=Stealth]
  \node[box] (intent) at (0,0) {Level\\intent};
  \node[box] (planner) at (3.3,0) {Local\\planner};
  \node[box] (turn) at (6.6,0) {Turning\\limiter};
  \node[box] (ctrl) at (9.9,0) {Velocity\\controller};
  \node[box] (motor) at (13.2,0) {Wheel\\effort};
  \draw[->] (intent) -- node[lbl, above] {$(v_i,\omega_i)$} (planner);
  \draw[->] (planner) -- node[lbl, above] {$(v,\omega)$} (turn);
  \draw[->] (turn) -- node[lbl, above] {$(v^\ast,\omega^\ast)$} (ctrl);
  \draw[->] (ctrl) -- node[lbl, above] {$u_L, u_R$} (motor);
  \draw[->] (intent.south) to[out=-30, in=-150] node[lbl, below] {L0, L2: planner bypassed} (turn.south);
\end{tikzpicture}
\caption{Command pipeline. L0 and L2 skip the local planner; every managed command passes the turning limiter; every command passes the velocity controller and its stop conditions.}
\label{fig:pipeline}
\end{figure}

\subsection{L0: teleoperation}

The intent is the operator's command, limited and checked for freshness:
\begin{equation}
  (v_i, \omega_i) = \begin{cases} \left(\operatorname{clip}(v_{\text{op}}, \pm 2),\; \operatorname{clip}(\omega_{\text{op}}, \pm 2)\right) & t - t_{\text{op}} < 0.5~\text{s}, \\ (0, 0) & \text{otherwise.} \end{cases}
\end{equation}
The command bypasses the planner: $(v, \omega) = (v_i, \omega_i)$.
From ARGOS it is a managed command, so it still passes the turning limiter and the velocity controller.
The operator closes the outer loop, and the 0.5~s lease means a lost link stops the rover within half a second.

\subsection{L1: assisted teleoperation}

The intent is the same operator command, but the planner chooses the motion:
\begin{equation}
  (v, \omega) = \arg\min_{(v,\omega) \in \mathcal{A}} J(v, \omega \mid v_{\text{op}}, \omega_{\text{op}}),
\end{equation}
where $\mathcal{A}$ is the admissible set.
On an open map the result equals the operator's command up to the acceleration window; near an obstacle the planner slows or steers away, and the decision is labelled \texttt{assisted}.

\subsection{Pursuit toward a goal}

L2 and L3 both use a pure-pursuit-style heading law toward a target point $\mathbf{g}$.
With $d = \lVert \mathbf{g} - (x, y) \rVert$ and bearing error $e_\theta = \operatorname{wrap}(\operatorname{atan2}(g_y - y,\, g_x - x) - \theta)$,
\begin{align}
  \omega_i &= \operatorname{clip}(k_\theta e_\theta,\; \pm\omega_{\max}^{\text{wp}}), \\
  v_i &= \begin{cases} v_c \left(1 - \dfrac{|e_\theta|}{\gamma}\right) \max\!\left(\min\!\left(\dfrac{d}{R_s}, 1\right),\; 0.2\right) & |e_\theta| < \gamma, \\ 0 & |e_\theta| \ge \gamma. \end{cases}
  \label{eq:pursuit}
\end{align}
The rover drives at cruise speed $v_c$ when aligned and far, slows linearly with heading error and stops forward motion beyond the gate $\gamma$, and slows to at least 20\% of cruise inside the slow radius $R_s$.
Inside the arrival radius $R_a$ it stops, or, if the goal has a final yaw $\theta_g$, turns with $\omega_i = \operatorname{clip}(k_\theta \operatorname{wrap}(\theta_g - \theta), \pm\omega_{\max}^{\text{wp}})$ until within 0.12~rad.
The defaults are $v_c = 1.0$~m/s, $k_\theta = 1.6$~s$^{-1}$, $\omega_{\max}^{\text{wp}} = 1.2$~rad/s, $\gamma = 0.55$~rad, $R_s = 3.0$~m, and $R_a = 0.75$~m.

\subsection{L2: direct waypoint}

L2 applies Equation~\eqref{eq:pursuit} straight toward the operator's goal and bypasses the planner, so it does not avoid obstacles.
It is the baseline that isolates the value of obstacle avoidance in L3.

\subsection{L3: waypoint with obstacle avoidance}

L3 adds a global route and the local planner.
An A* search on the occupancy grid finds a route of free cells from the rover to the goal.
The graph is four-connected with unit edge cost, which prevents cutting corners diagonally past obstacles, and a cell is traversable only if the footprint test of Section~\ref{sec:auto-planner} passes there.
The heuristic is the Manhattan distance in cells,
\begin{equation}
  h(c) = |c_{\text{col}} - g_{\text{col}}| + |c_{\text{row}} - g_{\text{row}}|,
\end{equation}
which is admissible and consistent for unit-cost four-connected motion, so the route is shortest in cell count.
The route is rechecked every 0.2~s and whenever the map changes along it.
The rover pursues the next route point with Equation~\eqref{eq:pursuit}, with the arrival radius tightened to 0.05~m at intermediate corners, and the planner turns that intent into the executed command.
Unknown space is never routed through.

\subsection{L4: frontier exploration}

L4 Explore drives to frontiers, free cells next to unknown space, with no per-goal approval, for a time budget $T_b$ set by the operator.
A breadth-first search from the rover over the grid visits each reachable cell $c$ with travel distance $n(c)$ in cells.
Its information gain $G(c)$ is the number of unknown cells in a square window of half-width $\lceil \rho / r \rceil + 2$ cells around it.
The chosen frontier maximizes gain per unit travel,
\begin{equation}
  f^\ast = \arg\max_{c \in \mathcal{F}} \frac{G(c)}{1 + n(c)},
  \label{eq:frontier}
\end{equation}
over the candidate set $\mathcal{F}$ of cells that pass the footprint test with no unknown-space exception, have $G(c) > 0$, lie more than $\max(\rho, 0.8~\text{m})$ from the rover, and are more than 1~m from any blacklisted point.
The frontier becomes a goal followed through the planner, with the horizon reduced to $H = 1$~s so that goals at the edge of the map remain reachable while unknown cells stay blocked.
If the rover makes no progress for 5~s, the frontier is blacklisted and another is chosen; after three attempts a point is excluded for the rest of the run.
Progress is reported as observed free area and explored fraction,
\begin{equation}
  A_{\text{free}} = N_{\text{free}}\,r^2, \qquad \eta = \frac{N_{\text{known}}}{N_{\text{cells}}}.
\end{equation}
The run ends when $t \ge T_b$, when $\mathcal{F}$ is empty, on a stop or takeover, or on a sustained stale map or safety hold.
The map-stale ending has a 3~s grace, while motion is still held from the first stale tick (Terra PR~\#60).
The explorer can also ask the arbiter for a look-around, a slow turn on the spot to observe more of the map.
That look-around was cleared only inside L4, so after an L4 run the rover kept executing it in L2 and L0 and reported \texttt{obstacle\_blocked} although those levels never consult the map; changing level now ends it (Terra PR~\#75).
In ARGOS this level is offered as ``L4 Explore'' with a 3, 5, or 10~minute budget (ARGOS PR~\#64).

\subsection{L4: target search}

Target search uses frontier navigation inside an operator-defined rectangle, implemented by marking every cell outside the rectangle as occupied, and confirms a target from detector observations $\{(x_k, y_k, c_k, t_k)\}$.
A target is confirmed when at least $N$ consecutive fresh observations of the requested class satisfy
\begin{equation}
  c_k \ge c_{\min}, \qquad \max_{j,k} \lVert (x_j, y_j) - (x_k, y_k) \rVert \le d_{\text{agree}}, \qquad t_{\text{last}} - t_{\text{first}} \ge \tau_{\text{span}},
\end{equation}
with no gap or age above $\tau_{\text{fresh}}$; a disagreeing observation restarts the set.
The defaults are $N = 3$, $c_{\min} = 0.8$, $d_{\text{agree}} = 1$~m, $\tau_{\text{span}} = 0.5$~s, and $\tau_{\text{fresh}} = 0.5$~s.
These are software thresholds, not a calibrated detector model.

\subsection{L5: task autonomy}
\label{sec:auto-l5}

L5 was defined as the level at which the robot decides what is needed.
Its first implementation is narrower: one rover takes a task in plain language from the operator and decides its own waypoints for it (Terra PR~\#70, ARGOS PR~\#85, merged October 25, 2026).
A language model running on the phone makes the decisions; the Terra core keeps every safety and motion rule.
\textbf{This level has not been run on a phone.}
At the time of writing TerraPhone with the model runner had not been built and run on an iOS~26 device, the ARGOS task panel had not been exercised against a live rover, and no task has been completed end to end.
Everything below describes the implemented design, not observed behaviour.

\paragraph{Loop.}
The operator selects L5 in ARGOS and types a task of at most 500 characters, with a time budget (default 300~s, at most 3600~s) and a limit on model calls (default 20, at most 100).
The arbiter then repeats one step until the task ends:
\begin{enumerate}
\item It builds a prompt from the occupancy map: the whole local map reduced to a symbolic grid of 1~m cells (a cell is occupied if any of its 0.1~m cells is occupied, free only if all of them are observed and free, and unknown otherwise), up to six reachable frontiers, up to eight obstacle blobs and eight landmarks, each with an identifier, the rover's pose, and the six most recent steps of history.
\item The phone runs the model on that prompt. The rover holds while the model thinks; the model runs off the control queue, so the 10~ms tick and the 20~Hz motor output do not wait for it.
\item The arbiter parses the reply, which must be one JSON object choosing one of four decisions: \texttt{go} (up to five waypoints, given as identifiers or local points), \texttt{ask\_operator} (a question shown in ARGOS; the rover waits for the answer), \texttt{done}, or \texttt{fail}.
\item Each waypoint is snapped to observed, footprint-clear free space within 0.5~m, and must have an A* route (Section~\ref{sec:auto-planner}) before the waypoint follower accepts it. Accepted waypoints are driven exactly as in L3.
\end{enumerate}
A reply that contains a speed, an effort or a servo command is rejected: the model chooses \emph{where}, never \emph{how}.
A waypoint whose route stays blocked for 4~s, or that is not reached within 90~s, is abandoned and the model is asked again.
A query with no answer after 90~s counts as a failed step, and three consecutive failed steps (invalid replies, timeouts or model errors) fail the task.
The operator can answer a question or add guidance at any time, which replaces the current plan, and can cancel; a stop, a level change or a goal cancel also ends the task.

\paragraph{Model.}
The model runner uses Apple's on-device Foundation Models framework (iOS~26 or later), with text input only.
The level is offered to ARGOS only while the phone reports that a model is available.
The runner sits behind a small protocol so that another on-device model, for example a vision--language model run through MLX or llama.cpp, can replace it.
\todo{Run L5 end to end on an iOS~26 iPhone, first in Zorvane and then on CrabTrapBot, and report model latency per step, the fraction of replies rejected and why, and task completion for a small set of fixed tasks.}

\paragraph{Relation to the rest of the ladder.}
L5 adds a decision maker above L3, not a new controller: the A* check, the local planner of Section~\ref{sec:auto-planner}, the turning limiter and every safety hold apply unchanged.
Because the model only proposes waypoints, a bad reply costs time, not safety.
The pilot study does not use L5.
Deciding what the \emph{fleet} needs, the original meaning of the level, is not implemented; Chapter~6 discusses the recommendation system that would sit between the operator and such decisions.

\begin{table}[ht]
\centering
\small
\caption{How each level produces motion. ``Managed'' commands pass the turning limiter.}
\label{tab:level-math}
\begin{tabularx}{\textwidth}{@{}l >{\raggedright\arraybackslash}X >{\raggedright\arraybackslash}p{2.2cm} >{\raggedright\arraybackslash}p{2.4cm}@{}}
\toprule
Level & Intent $(v_i, \omega_i)$ & Local planner & Human role \\
\midrule
L0 Teleop & Operator command, 0.5~s lease & Bypassed & Closes the loop \\
L1 Assisted & Operator command & Eq.~\eqref{eq:planner-cost} & Drives; planner vetoes \\
L2 Direct & Pursuit, Eq.~\eqref{eq:pursuit}, to goal & Bypassed & Sets goal; watches \\
L3 Waypoint & Pursuit along A* route & Eq.~\eqref{eq:planner-cost} & Sets goal \\
L4 Explore & Pursuit to $f^\ast$, Eq.~\eqref{eq:frontier} & Eq.~\eqref{eq:planner-cost}, $H = 1$~s & Sets budget \\
L4 Search & Pursuit to frontier in bounds; confirmation rule & Eq.~\eqref{eq:planner-cost} & Sets area, class, budget \\
L5 Task & Pursuit along A* route to model-chosen waypoints & Eq.~\eqref{eq:planner-cost} & Gives task; answers questions \\
\bottomrule
\end{tabularx}
\end{table}

\section{Outdoor guidance and heading fusion}
\label{sec:auto-outdoor}

\subsection{Geographic goals}

Outdoors, the goal is a latitude and longitude.
From the current fix $(\varphi_1, \lambda_1)$ to the goal $(\varphi_2, \lambda_2)$, Terra computes the great-circle distance and initial bearing with the haversine formula, using $R_\oplus = 6\,371\,000$~m:
\begin{align}
  h &= \sin^2\!\frac{\Delta\varphi}{2} + \cos\varphi_1 \cos\varphi_2 \sin^2\!\frac{\Delta\lambda}{2}, \qquad d = 2R_\oplus \arcsin\sqrt{h}, \\
  \beta &= \operatorname{atan2}\!\left(\sin\Delta\lambda \cos\varphi_2,\; \cos\varphi_1 \sin\varphi_2 - \sin\varphi_1 \cos\varphi_2 \cos\Delta\lambda\right).
\end{align}
The bearing $\beta$ (degrees clockwise from north) is converted into a local yaw $\psi$ using the fused heading of the next subsection, and a local target is placed at most 3~m ahead:
\begin{equation}
  \mathbf{g} = (x, y) + \min(d, 3)\,(\cos\psi, \sin\psi).
\end{equation}
That target is refreshed at most twice per second and followed exactly as in L3, so the local planner still checks every command; forward speed is capped at 0.3~m/s.
Arrival requires three distinct fixes within the arrival radius.
For goals given in a local frame with a geographic origin $(\varphi_0, \lambda_0)$, Terra uses an equirectangular projection, $x_{\text{east}} = (\lambda - \lambda_0)\,k\cos\varphi_0$ and $y_{\text{north}} = (\varphi - \varphi_0)\,k$, with $k = 40\,075\,016.686 / 360$~m per degree.

\subsection{Heading fusion}

GPS gives position, not heading, and the phone's compass near the motors was off by 30--50$^\circ$.
Terra therefore estimates a single slowly varying offset $c$ (degrees) between the local VIO frame and true north, such that the true heading is
\begin{equation}
  h = c - \theta^{\circ},
\end{equation}
where $\theta^{\circ}$ is the local yaw in degrees (local yaw is counter-clockwise; true heading is clockwise).
VIO carries the heading between corrections, and the filter only has to track $c$.
It is a scalar Kalman filter with random-walk process noise.
Between corrections the variance grows as
\begin{equation}
  P(t) = P_k + q\,(t - t_k), \qquad q = 0.25~\text{deg}^2/\text{s}.
\end{equation}
Two measurements of $c$ are available.
A compass sample $h_c$ with reported accuracy $\sigma_c$ gives $z = h_c + \theta^{\circ}$ with $R = \sigma_c^2$.
A GPS course-over-ground sample $\chi$ gives $z = \chi + \operatorname{atan2}(\Delta y, \Delta x)^{\circ}$, pairing it with the local direction of travel over the same 1~s window, with $R = \max(\sigma_\chi, 2^\circ)^2$; it is used only at 0.5~m/s or faster, with less than 10$^\circ$ of turning in the window.
For a measurement $z$, the innovation and update are
\begin{align}
  \nu &= \operatorname{wrap}_{180}(z - c), \\
  K &= \frac{P}{P + R}, \qquad c \leftarrow c + K\nu, \\
  P &\leftarrow \begin{cases} \dfrac{PR}{P + R} & \text{GPS course,} \\[1ex] \min\!\left(P,\; \max\!\left(\dfrac{PR}{P + R},\; R\right)\right) & \text{compass.} \end{cases}
\end{align}
The compass update floors the variance at $R$, because compass error is mostly a bias that does not average out; a compass sample is ignored entirely when $R > 4P$.
A measurement is rejected if $\nu^2 > 9(P + R)$, a three-sigma gate; after ten consecutive rejections the filter restarts from the latest measurement.
The filter also restarts if the VIO pose has a gap longer than 2~s or a yaw step faster than 3~rad/s.
The reported heading accuracy is $\sqrt{P(t)}$; ARGOS holds an outdoor mission when it exceeds 35$^\circ$, and Terra requires it below 30$^\circ$ to admit one.

\section{Multi-rover map fusion (prototype)}
\label{sec:auto-fusion}

Every rover in Section~\ref{sec:auto-map} maps in its own frame, and nothing on the rovers combines those maps.
A first step towards a shared map is a two-rover prototype, \texttt{scripts/plot\_sim\_rovers.py} in the Terra repository, which registers one rover's occupancy grid onto another's and shows the fused result.
\todo{The script is not yet committed to Terra (checked Oct 26, 2026); this section was checked against the working copy Yojan shared that day. Commit it and pin the revision in the header comment.}
It is a visualisation and registration tool only.
It runs as a separate Python process that subscribes over Zenoh, it does not publish the fused grid back, and it does not change either rover's mapper or the arbiter.
A fielded system shows what a full version needs.
In Team CERBERUS's SubT system, each robot sent submaps to a central server at the base station whenever it had a link, and the server closed visual and LiDAR loops across robots and re-optimized all trajectories together~\citep{tranzatto2022finals}.
A central coordinator then used the merged map to find global frontiers and send robots to them, while each robot's onboard planner kept exploring when the link was lost~\citep{kulkarni2022teamed}.
The prototype here covers only the first part of that pipeline, registration and display.

\subsection{Inputs and accumulation}

The prototype subscribes to the \texttt{pose} and \texttt{map/occupancy} topics of a reference rover $R$ and a source rover $S$ (by default IDs 0 and 1), and estimates the rigid transform that carries $S$'s map into $R$'s frame.
The occupancy packet is the one Terra already publishes (Section~\ref{sec:auto-map}): values $-1$ for unknown and $0$--$100$ otherwise, rows along world $+y$, origin at the minimum cell corner, and the same threshold $P \ge 65$ for occupied.
The pose is drawn on the plot but is not a registration measurement; the transform comes from map structure alone.

Each published grid is a rolling window, so cells fall out of it as the rover moves.
The prototype keeps a persistent map per rover, keyed by integer lattice coordinates.
A cell at column $j$ and row $k$ of a grid with origin $(o_x, o_y)$ and resolution $r$ has key
\begin{equation}
  (i_x, i_y) = \left(\operatorname{round}(o_x / r) + j,\; \operatorname{round}(o_y / r) + k\right).
\end{equation}
This is exact for Terra's mapper, which only ever shifts its origin by whole cells.
The newest known value at a key replaces the old one, unknown cells never erase evidence, and older or repeated map revisions are ignored.
A change of run ID clears both maps and the current alignment.
When a persistent map would exceed $5 \times 10^5$ cells it is coarsened by an integer factor $s = \max(2, \lceil\sqrt{N_x N_y / 5 \times 10^5}\,\rceil)$, keeping the maximum value in each bin, which favours obstacles.

\subsection{Registration}

Registration uses only occupied cells.
Their centres form point sets $P_R$ and $P_S$.
Each set is downsampled to one point per voxel of width $\max(0.25~\text{m}, 2r)$ and capped at 250 points, so that long walls do not dominate.

\paragraph{Global search.}
Before an alignment exists, and then every 30~s, the prototype sweeps yaw $\theta$ in $10^\circ$ steps over the full circle.
At each $\theta$, every pair of a reference point and a rotated source point votes for a translation,
\begin{equation}
  \mathbf{d}_{ij} = \mathbf{p}_{R,i} - R(\theta)\,\mathbf{p}_{S,j},
\end{equation}
quantised into bins of width $b = \max(0.25~\text{m}, 3r)$.
The two fullest bins at each yaw become candidates.
A candidate is scored by nearest-neighbour matching within $\rho = \max(0.35~\text{m}, 4r)$, with each reference point matched at most once.
With $M$ unique matches, $B$ occupied 1~m bins containing match midpoints, and RMS match distance $e$, candidates are ranked by larger $M$, then larger $B$, then smaller $e$, then more votes.
This is an exhaustive discrete search, not RANSAC.

\paragraph{Refinement.}
The best candidate, or between global searches the current alignment, seeds point-to-point ICP on up to 6\,000 occupied points per map.
Each iteration pairs every transformed source point with its nearest reference point, drops pairs farther than $\max(0.65~\text{m}, 6r)$, trims the rest to the 75th percentile of distance, and requires at least eight pairs.
The rotation is the least-squares solution from the SVD of
\begin{equation}
  H = \sum_i (\mathbf{x}_i - \bar{\mathbf{x}})(\mathbf{y}_i - \bar{\mathbf{y}})^{\mathsf{T}},
\end{equation}
and the translation aligns the centroids.
ICP stops after 35 iterations or when the update falls below $10^{-4}$~m and $10^{-5}$~rad.
The refined transform is then rescored on the downsampled sets at radius $\rho$.

\subsection{Scoring and acceptance}

With precision $P_m = M / N_S$ and recall $Q_m = M / N_R$ over the downsampled point counts, the prototype computes a confidence
\begin{equation}
  C = \operatorname{clip}_{[0,1]}\!\left(0.55\,F_1 + 0.20\min(1, B/8) + 0.20\,e^{-e/\rho} + 0.05\,A_K\right),
  \qquad F_1 = \frac{2 P_m Q_m}{P_m + Q_m},
\end{equation}
where $A_K$ is the mean agreement $1 - |c_R - c_S|/100$ over cells known in both maps, computed only during a global search.
It also computes an evidence score
\begin{equation}
  I = \operatorname{clip}_{[0,1]}\!\left(0.40\min(1, \tfrac{M}{40}) + 0.25\min(1, \tfrac{B}{8}) + 0.20\min(1, \tfrac{\min(N_R, N_S)}{100}) + 0.15\min(1, \tfrac{N_{\text{ICP}}}{50})\right),
\end{equation}
and their product $J = CI$.
The weights are hand-chosen.
$C$ is a ranking score, not a calibrated probability that the alignment is correct, and the thesis does not use it as one.

A candidate is eligible if $M \ge 6$, $B \ge 2$, $e \le 0.45$~m, $C \ge 0.18$ and $I \ge 0.025$.
The first eligible candidate becomes the incumbent.
A later one replaces it only if $C$ rises by more than $10^{-4}$, $I$ does not fall, and $J$ rises.
A rejected candidate leaves the current fused map in place.
Alignment is attempted at most every 2~s.

\subsection{Fusion}

The fused grid is recomputed from the two persistent maps under the incumbent transform, never from the previous fused grid, so no observation is counted twice.
It covers the union of both extents at the finer input resolution, coarsened if needed to stay within $5 \times 10^5$ cells.
Each known value $c$ becomes $p = \operatorname{clip}(c/100, 0.01, 0.99)$, and the two maps are combined in log odds,
\begin{equation}
  \ell_F = \operatorname{clip}\!\left(\sum_{m \in \{R, S\}} \ln\frac{p_m}{1 - p_m},\; \pm 4\right),
  \qquad P_F = \operatorname{round}\!\left(\frac{100}{1 + e^{-\ell_F}}\right),
\end{equation}
mirroring the single-rover update and its clamp.
A cell unknown in both maps stays $-1$.
The sum treats the two rovers' observations of a cell as independent evidence.
Fusion runs at most every 0.75~s and the plot at most at 2~Hz.

\subsection{Validation so far}

In Zorvane all rovers report maps in one shared world frame (Chapter~\ref{ch:simulation}), so the correct transform between two simulated rovers is the identity.
To test registration on simulated data the prototype therefore applies a known transform $A = (2.0~\text{m}, -1.5~\text{m}, 23^\circ)$ to the source map by default and reports the residual between the estimate and $A^{-1}$; \texttt{--no-test-transform} disables it for maps that are genuinely in different frames, as two phones' VIO frames are.
A built-in self-test registers a synthetic asymmetric $24 \times 20$~m layout at 0.1~m under three known transforms and passes if each alignment is eligible, each residual is within 0.30~m and $4^\circ$, and the fused map keeps at least 70\% of the reference map's known cells.
Table~\ref{tab:fusion-selftest} gives the result of one run.

\begin{table}[ht]
\centering
\begin{tabular}{rrrrrr}
\toprule
\multicolumn{3}{c}{Applied transform} & \multicolumn{2}{c}{Residual} & \\
\cmidrule(lr){1-3}\cmidrule(lr){4-5}
$x$ (m) & $y$ (m) & yaw & translation (m) & yaw & $M$ \\
\midrule
2.0 & $-1.5$ & $23^\circ$ & 0.028 & $0.97^\circ$ & 160 \\
$-3.0$ & 2.0 & $-31^\circ$ & 0.035 & $0.60^\circ$ & 159 \\
0.8 & 1.2 & $0^\circ$ & 0.199 & $0.01^\circ$ & 170 \\
\bottomrule
\end{tabular}
\caption{Prototype self-test on synthetic data (run Oct 26, 2026). $M$ is the number of unique occupied-cell matches after ICP.}
\label{tab:fusion-selftest}
\end{table}

The rotated cases recover the transform to a few centimetres.
The pure translation, which should be the easiest, is the worst: ICP converges two cells short along $x$ with an RMS match distance of only 0.01~m, because points on long axis-aligned walls match exactly whatever their position along the wall, and the 75th-percentile trim discards the wall ends that would correct it.
It still passes, but with only 0.10~m to spare.
The coverage check passes for any transform, since the reference map alone supplies its known cells, so it guards against coarsening losses and not against a wrong merge.
This is a software check on synthetic data, not evidence of performance on real maps.
\todo{Report registration residuals from the Zorvane two-rover run with the test transform, and from two phones in one room with \texttt{--no-test-transform} against OptiTrack or tape-measured ground truth.}

\subsection{Limitations}

\begin{itemize}
  \item \textbf{Uncalibrated scores.} $C$, $I$ and $J$ rank candidates; none is a probability of a correct merge.
  \item \textbf{Frame stability.} Accumulation assumes each rover's frame does not jump. An ARKit relocalisation or VIO reset without a new run ID would corrupt that rover's persistent map.
  \item \textbf{Ratcheting incumbent.} A new alignment, local or global, replaces the incumbent only if its confidence is strictly higher. A wrong early merge with a high score is therefore kept until the run changes, and since the incumbent was scored on smaller maps, a correct transform rescored on grown maps can be rejected because recall over the larger point sets fell.
  \item \textbf{Static world.} Old observations persist until the cell is seen again, and max-coarsening keeps stale obstacles.
  \item \textbf{Ambiguity.} Translation voting can lock onto a wrong but plausible alignment in symmetric or repetitive spaces such as corridors; the spatial-spread term reduces this but does not remove it.
  \item \textbf{Two rovers.} It registers one source to one reference. It is not a multi-robot pose graph and has no loop closure.
  \item \textbf{Not on the robots.} The fused map reaches no rover, planner, or ARGOS. Exploration still runs on each rover's own map, so two rovers can explore the same area.
\end{itemize}

Publishing the fused map back, so that the frontier explorer of Section~\ref{sec:auto-levels} can avoid space another rover has already covered, is future work.

\section{Summary}

Each autonomy level is a different source of intent feeding one shared pipeline: a dynamic-window planner that only accepts commands it can still stop from, a turning limiter that keeps skid-steer turns gentle, and a feedforward PI velocity controller on VIO and IMU estimates, with every stage stopping on stale input.
L0 and L2 bypass the planner, which makes them the baselines for measuring what obstacle avoidance adds; L1, L3, L4, and L5 share it.
L5 puts an on-phone language model above L3 to choose waypoints for a plain-language task, with every waypoint checked by A* and driven on the L3 path; it is implemented but untested on hardware.
The occupancy map's ground plane is level and placed at the mount-calibrated camera height, a correction found by inspecting labelled point clouds streamed to ARGOS, and ARKit's classified floor mesh now vetoes obstacle returns that lie on the floor; wall, window and door mesh now stops rays at glass, though this is untested on real glass.
Against motion capture, the phone's VIO speed was steady, low-noise and correctly oriented but read about 9\% low in the capture room, which the speed loop tolerates and absolute distances do not.
The faults that most affected driving in October 2026 were timing faults on the phone's control queue, not errors in the control laws: each stale input correctly produced a hold, and the holds compounded.
Outdoors, a scalar Kalman filter on the VIO-to-north offset fuses compass and GPS course into the heading that turns a geographic bearing into a local goal.
A separate two-rover prototype registers one rover's occupancy grid onto another's by occupied-cell voting and ICP and fuses them in log odds, but only for display; no rover uses the fused map yet.
\todo{Add a figure of a recorded planner trace (candidates and selected arc) from a real run.}
\todo{Report step responses of the velocity controller on the real rover once it is tuned.}
